\documentclass[11pt]{article}
\usepackage[margin=1in]{geometry}
\usepackage{amsmath,amssymb,amsthm,bm,xcolor,hyperref,booktabs,enumitem}
\newtheorem{proposition}{Proposition}
\newtheorem{assumption}{Assumption}
\newtheorem{definition}{Definition}
\newtheorem{theorem}{Theorem}
\newtheorem{example}{Example}

\newcommand{\TESTED}[1]{{\color{teal}\small$\blacktriangleright$\,\textbf{[EVIDENCE]} #1}}
\newcommand{\PARTIAL}[1]{{\color{orange}\small$\blacktriangleright$\,\textbf{[PARTIAL]} #1}}
\newcommand{\PROP}[1]{{\color{violet}\small$\blacktriangleright$\,\textbf{[PROPOSED / UNTESTED]} #1}}
\newenvironment{intuition}{\par\smallskip\noindent\hangindent=1.5em\textbf{\itshape Intuition.}\itshape\ \ignorespaces}{\par\smallskip}

\newcommand{\R}{\mathbb{R}}
\newcommand{\E}{\mathbb{E}}
\newcommand{\N}{\mathcal{N}}
\renewcommand{\Pr}{\mathbb{P}}

\title{\textbf{Bayesian Inference under an Imperfect LLM Oracle:\\ A Structured Measurement-Error Model for
Relevance, Identified by Corpus Structure}\\[4pt]
\large\normalfont(expanded, hand-holding draft --- Statistics paper of a two-paper split)}
\author{Robert Richardson}
\date{\today}

\begin{document}
\maketitle

\begin{abstract}
Large language models are increasingly used as cheap ``oracles''---relevance judges, correctness verifiers,
data annotators---and their outputs are almost always treated either as ground truth or as exchangeably noisy
labels. We argue this is a measurement problem in disguise, and a badly specified one: LLM judgments carry
\emph{systematic, structured} bias, not random noise. In multi-hop relevance judging we document a concrete
instance---\emph{bridge-blindness}, where a passage essential to a reasoning chain is judged irrelevant
\emph{because} it does not directly answer the question---and show the bias is a function of a latent passage
\emph{role}, violating the conditional-random-error assumption of classical rater models. We develop a Bayesian
measurement-error model in which the ordinal LLM judgment is a noisy view of a latent role with a role-dependent
confusion law, and---the central methodological contribution---we use the \emph{corpus graph} to identify the
model \emph{without} gold labels or repeated annotation: relevance is smooth on the graph, so a confidently
relevant anchor and its neighbours provide correlated latent roles viewed through the same instrument,
separating the judge's bias from the truth. We state identification conditions, give an inference procedure,
and target calibrated relevance posteriors whose correction demonstrably improves a downstream decision. The
application is passage retrieval, but the object---rigorous inference under a biased LLM oracle---recurs across
weak supervision, LLM annotation, and automated evaluation.
\end{abstract}

\paragraph{How to read this draft.} Almost everything here is a \emph{proposal}; this document's job is to make
the target precise and confirm the fundamentals. Grey ``Intuition'' notes give the plain-language version.
Status tags: \TESTED{} supported by measurements we already have;\ \PARTIAL{} partial;\ \PROP{} proposed, not
yet run.

\tableofcontents

\section{Introduction}
\subsection{The setting}
An \emph{LLM-as-judge} is a large language model prompted to score something---here, ``given query $q$ and
passage $d$, how relevant is $d$?''---returning an ordinal grade $g\in\{0,1,2\}$. It is cheap and scalable, so
pipelines lean on it heavily: to filter retrieved passages, to label training data, to evaluate other systems.
The near-universal practice is to \emph{believe} the judge: threshold $g$ and treat the result as the truth,
or average several judges as if their disagreements were random noise.

We find both practices are wrong in the same way. A relevance judge, however carefully prompted, systematically
under-rates \emph{bridge} passages---those that supply an intermediate linking fact rather than the final
answer. Ask ``is this passage relevant to \emph{who is the spouse of the performer of Green}?'' about the
passage ``$X$ performed \emph{Green}'' and the judge often says \emph{no}: the passage does not mention spouses.
Yet that passage is essential---it names $X$, without whom the chain cannot be completed. The judge's error is
not random; it is a deterministic function of the passage's \emph{role} in the reasoning.

\subsection{Why this is a statistics problem, and a nonstandard one}
Observing a biased, noisy version of a latent truth is the province of \emph{measurement-error} and
\emph{misclassification} models. But this instance has two features that break the standard toolkit:
\begin{enumerate}[leftmargin=*,itemsep=1pt]
  \item \textbf{The misclassification is \emph{differential} with respect to the scientific target.} Our target
  is binary relevance $r_i$. Two truly-relevant passages ($r_i{=}1$) have radically different grade
  distributions depending on an unobserved functional role---bridge vs.\ direct---so relevance alone does not
  render the measurement non-differential; the error law is fixed only after \emph{refining} $r_i$ by a latent
  role $z_i$. This is the feature a standard binary-relevance rater model cannot represent, and it is the sense
  in which the problem is nonstandard---not that role-dependence per se defeats latent-class models (it does
  not; see \S\ref{sec:prelim}).
  \item \textbf{We usually have one judge and no gold.} The classical routes to pinning down the error
  mechanism---many independent raters per item, or a labeled subset---are often unavailable at scale, which is
  why identification (\S\ref{sec:ident}) is the crux.
\end{enumerate}
Our contribution is a model for (1) and a use of \emph{corpus structure} to solve (2). The payoff is a
calibrated posterior over true relevance computed from biased judgments, and a transportable template for how a
statistician should consume an LLM oracle.

\subsection{Relation to the companion paper}
The companion CS/ML paper uses the corpus graph as the \emph{covariance} of a latent relevance field: structure
propagates \emph{belief} (judge one passage, update its neighbours). Here the \emph{same} graph is the
\emph{instrument that identifies the measurement model}: structure propagates \emph{calibration} (a known-relevant
anchor reveals how the judge behaves on its neighbours' roles). One object, two epistemic jobs. The two papers
are disjoint in contribution---a retrieval method vs.\ a measurement model---and cite each other.

\section{Background and Preliminaries}\label{sec:prelim}
\subsection{Measurement error and misclassification}
\emph{Measurement error} is the study of inference when a variable of interest is observed only through a noisy
proxy. For continuous data the archetype is $W=X+U$ (observe $W$, want $X$); naively using $W$ for $X$ biases
downstream estimates, and correcting requires a model of $U$. For \emph{categorical} data the analogue is
\emph{misclassification}: the observed category is a corrupted version of the true one, described by a
\emph{confusion} (or misclassification) matrix $\Pi$ with $\Pi_{c\ell}=\Pr(\text{observe }\ell\mid\text{true
}c)$, rows summing to one. The lesson is the same: ignoring $\Pi$ biases everything; correcting requires
knowing (or inferring) $\Pi$.

\subsection{Latent-class rater models (Dawid--Skene)}
The canonical model for \emph{learning $\Pi$ and the truth together} is Dawid--Skene. Each item $i$ has an
unknown true label $z_i$ drawn from a class prior; each rater $\rho$ has a confusion $\Pi^{(\rho)}$; the observed
labels $g_i^{(\rho)}$ are conditionally independent across raters given $z_i$, each drawn from row $z_i$ of
$\Pi^{(\rho)}$. One fits $\{z_i\},\{\Pi^{(\rho)}\}$ by EM or Bayesian inference. Two assumptions do the work:
(a) \emph{conditional independence} of raters given the truth, and (b) errors are \emph{random} given the true
label (the confusion depends on nothing but $z_i$). Identification (below) relies on \emph{replication}: several
raters per item.
\begin{intuition}
Dawid--Skene says: if three annotators look at the same item and two say ``cat,'' the item is probably a cat and
the third annotator is probably noisy---and by watching many items you learn \emph{how} noisy each annotator is.
It is the workhorse of crowdsourced labeling. Our problem does \emph{not} break latent-class rater models per
se: once the latent \emph{role} $z_i$ is the class, $g_i\mid z_i$ is a perfectly standard class-dependent error
law. The sharper statistical point is that this misclassification is \emph{differential} with respect to our
target, binary relevance $r_i$: two $r_i{=}1$ passages of different roles have different grade laws, so $r_i{=}1$
is not enough to fix the error law---you must refine $r_i$ by $z_i$. And we typically lack the replication that
Dawid--Skene exploits, which is what makes identification (\S\ref{sec:ident}) nontrivial.
\end{intuition}

\subsection{Ordinal responses and the probit/threshold parameterization}
An \emph{ordinal} response like $g\in\{0<1<2\}$ is naturally modeled by a latent continuous ``propensity''
$\eta$ crossing ordered thresholds $\kappa_1<\kappa_2$: $g=\sum_t\mathbf 1\{\eta>\kappa_t\}$. Covariates (here,
the latent role) shift $\eta$. This is the ordered-probit model; it is a parsimonious way to write a confusion
matrix that respects the ordering (a bridge is not mislabeled uniformly---it is pushed \emph{down} toward
$g{=}0$), with far fewer free parameters than an unconstrained $3\times3$ $\Pi$.

\subsection{Markov random fields on graphs: the Potts prior}
To say ``the truth is smooth on the corpus graph'' for \emph{discrete} labels we use a \emph{Potts model}, the
categorical cousin of the GMRF: $\Pr(z)\propto\exp\!\big(\theta\sum_{(i,j)\in A}\mathbf 1\{z_i\!=\!z_j\}\big)$,
which rewards graph-neighbours for sharing a label, with $\theta>0$ setting the coupling strength. We will use a
\emph{generalized} version that rewards \emph{compatible} roles (a chain is a connected direct--bridge
structure, not a monochrome blob).
\begin{intuition}
A Potts prior is a soft constraint: ``passages that are linked in the corpus tend to have related roles.'' It is
where we encode the belief that a reasoning chain is a \emph{connected} object, which is precisely what lets a
known anchor say something about its neighbours' hidden roles.
\end{intuition}

\subsection{Identifiability, and why it is the whole game here}\label{sec:idprelim}
A model is \emph{identified} if different parameter values imply different distributions of the observed data;
otherwise no amount of data can separate them. Confusion/latent-class models are notoriously \emph{non}-identified
without extra structure, chiefly through \emph{label switching}: you can permute the latent labels and
compensate by permuting the confusion rows, leaving the observed-label distribution unchanged. Concretely, from
the \emph{marginal} distribution of a single judge's grades you cannot separate ``this passage is irrelevant''
($z{=}0$) from ``this passage is a relevant bridge, and the judge is blind to bridges'' ($z{=}1$ with confusion
mass on $g{=}0$): both produce $g{=}0$. The standard escapes are (i) \emph{replication} (many conditionally
independent raters, à la Dawid--Skene), or (ii) an \emph{anchor}/gold subset with known labels. Our proposal
adds a third escape: (iii) \emph{structure}. This is the crux of the paper (\S\ref{sec:ident}).

\section{The Structured Measurement-Error Model}\label{sec:model}
\subsection{Latent roles}
Fix a query with candidate pool $\{d_1,\dots,d_N\}$. Each passage has a latent \emph{role}
\begin{equation}
  z_i\in\{\underbrace{0}_{\textsf{irrelevant}},\ \underbrace{1}_{\textsf{bridge}},\ \underbrace{2}_{\textsf{direct}}\},
  \qquad \text{true relevance } r_i=\mathbf 1\{z_i\ge1\}.
\end{equation}
Bridge and direct passages are \emph{both} relevant but functionally different: the direct passage contains (part
of) the answer; the bridge supplies a linking fact. This three-way latent is the modeling heart, because the
judge's bias lives at exactly the bridge/direct distinction, which a binary relevant/irrelevant latent cannot
express.

\begin{example}[running example]
For ``spouse of the performer of Green,'' the passage ``$X$ performed Green'' has $z=1$ (bridge) and ``$X$'s
spouse is $Y$'' has $z=2$ (direct). A good judge rates the direct passage highly ($g{=}2$) but frequently rates
the bridge $g{=}0$---the error we want to model and correct.
\end{example}

\subsection{The observation law (role-dependent confusion)}
A judge $\rho$ (a specific model + prompt) returns $g_i^{(\rho)}$ from a role-dependent confusion,
\begin{equation}\label{eq:obs}
  g_i^{(\rho)}\,\big|\,z_i=c \;\sim\; \mathrm{Categorical}\!\big(\pi^{(\rho)}_c\big),\qquad
  \pi^{(\rho)}_c\in\Delta^2,\quad \Pi^{(\rho)}=\big(\pi^{(\rho)}_0,\pi^{(\rho)}_1,\pi^{(\rho)}_2\big)^\top .
\end{equation}
\emph{Bridge-blindness} is the structural feature $\Pr(g{=}0\mid z{=}1)\gg\Pr(g{=}0\mid z{=}2)$: relevant
bridges are frequently graded irrelevant, relevant direct passages are not. Rather than leave $\Pi$ fully free,
we prefer the ordered-probit form, which builds the ordering in and isolates the bias in one interpretable
parameter per role,
\begin{equation}\label{eq:probit}
  g_i^{(\rho)}=\sum_{t\in\{1,2\}}\mathbf 1\{\eta_i^{(\rho)}>\kappa_t^{(\rho)}\},\qquad
  \eta_i^{(\rho)}=\alpha^{(\rho)}_{z_i}+\varepsilon_i,\ \ \varepsilon_i\sim\N(0,1),
\end{equation}
where $\alpha^{(\rho)}_c$ is the role's ``propensity to be graded up.'' Bridge-blindness is the statement
$\alpha^{(\rho)}_1$ is depressed toward $\alpha^{(\rho)}_0$ despite bridges being relevant---a single,
estimable, interpretable bias parameter.
\begin{intuition}
Equation \eqref{eq:probit} is just ``each passage has a hidden propensity to be rated well; its role sets that
propensity; noise does the rest.'' Bridge-blindness is one number: the gap $\alpha_2-\alpha_1$ that the judge
\emph{shouldn't} have (both are relevant) but does.
\end{intuition}
\TESTED{The role dependence is real. A binary answer-relevance judge has gold-recall $\approx0.35$; an
explicitly hop-aware prompt raises it to $0.82$---i.e., prompt framing shifts $\alpha_1$, exactly the mechanism
in \eqref{eq:probit}---yet precision stays low ($0.17$--$0.36$) and gold-recall degrades on longer chains and
adversarial corpora ($0.89\to0.57$). These are the sufficient statistics the model must reproduce.}

\subsection{The prior on latent roles via corpus structure}
We tie the roles to (a) the dense retriever and (b) the corpus graph $A$ through a generalized Potts prior,
\begin{equation}\label{eq:prior}
  \Pr(z)\;\propto\;\exp\!\Big(\underbrace{\sum_i\psi(z_i;\,m_i)}_{\text{semantic anchor}}\;+\;
  \underbrace{\theta\sum_{(i,j)\in A}\omega(z_i,z_j)}_{\text{graph coupling}}\Big),
\end{equation}
where $\psi(\cdot;m_i)$ pulls high-cosine passages toward relevance (the retriever prior $m_i$ from Paper A),
and $\omega$ rewards graph-connected passages for taking \emph{compatible} roles (e.g.\ a bridge next to a
direct passage). This is the \emph{identifying} prior: it encodes that the truth is graph-structured even
though the judge is graph-blind.

\section{Identification by Structure (the crux)}\label{sec:ident}
\subsection{The problem, restated}
With one judge and no gold, \eqref{eq:obs} alone does not identify $(\Pi,z)$: a passage graded $g{=}0$ could be
truly irrelevant ($z{=}0$) or a bridge the judge is blind to ($z{=}1$), and the marginal grade distribution
cannot tell them apart (\S\ref{sec:idprelim}). We need an outside constraint that behaves like replication.

\subsection{The right analogy is a hidden Markov model, not Dawid--Skene}
Dawid--Skene buys identification with \emph{replication}: several conditionally-independent raters view the
\emph{same} latent state. We have one rater, so that route is closed. The correct analogy is instead a
\emph{hidden Markov model} (HMM)---or, on a general corpus, a \emph{hidden Markov random field}. In an HMM each
latent state is observed only \emph{once}, yet the \emph{dependence among the latent states} induces dependence
among those single observations, and under rank/separation conditions the emission and transition laws are
identified \emph{up to a permutation of the states, with no replication}. Our reasoning chains are literally
such a process: the gold evidence for a $k$-hop question is a sequence of roles
$z_1\!\to\!z_2\!\to\!\cdots\!\to\!z_m$ tied by a role-transition structure, and the judge emits one ordinal
grade per passage. \emph{The chain's own transitions manufacture the multiple views a lone judge cannot supply.}

\subsection{Why this is not the clusterability route (heterophily is the point)}
A known and tempting shortcut identifies label noise from \emph{local same-class} structure: if graph-neighbours
share the true class, their single noisy labels furnish repeated views of one latent state, and higher-order
consensus identifies the confusion matrix without anchors (the \emph{clusterability} route, Zhu et al., 2021).
Our structure is \emph{not} homophilic in that sense. A direct-evidence passage adjacent to a bridge has a
\emph{different} role, $z_i\neq z_j$---but a \emph{compatible} one, which is exactly why the prior
\eqref{eq:prior} rewards role \emph{compatibility} rather than role \emph{identity}. Clusterability's consensus
argument therefore does not apply; the identifying object is instead the role \emph{transition/compatibility}
matrix $T$ itself, recovered through the triple law \eqref{eq:triple}. The novelty we can actually defend is
consequently narrow and specific---\emph{single-judge identification of a role-dependent measurement channel
from heterophilic, role-compatible corpus structure}---not the broad and long-established claim that
``dependence can substitute for replication'' (true since classical HMM identifiability; Gassiat et al.). A
bonus falls out: estimating $T$ recovers the latent \emph{grammar} of multi-hop evidence---which functional
roles follow which---an object clusterability's trivial same-class structure does not possess.

\subsection{A first identification theorem (the chain case)}
Take latent roles on a chain, $z_1\!\to\!\cdots\!\to\!z_m$, with a role-transition matrix $T$ and one-judge
ordinal emissions $g_i\mid z_i\sim p_\theta(g\mid z)$. Across many \emph{independent} query-chains we observe
the joint law of consecutive grades. The key object is the distribution of a \emph{triple},
\begin{equation}\label{eq:triple}
  \Pr(g_i,\,g_{i+1},\,g_{i+2}),
\end{equation}
which, by the Markov property, factorizes through the shared latent $z_{i+1}$ and so furnishes \emph{three}
conditionally-independent views of that one hidden state. This is exactly the configuration in which Kruskal's
theorem on three-way tensor decomposition delivers \emph{generic} identifiability of finite mixtures and HMMs
(Allman, Matias \& Rhodes, 2009): the emission $p_\theta(g\mid z)$ and the transition $T$ are recovered up to a
relabeling of roles, provided the role-emission distributions are linearly independent---a separation condition
our model satisfies whenever bridge and direct grades differ (precisely the bridge-blindness gap
$\alpha_2-\alpha_1$).

\PROP{\textbf{Theorem (chain identification --- to prove).} Under a chain role-process with linearly-independent
role-emission distributions and $m\ge3$, the emission law $p_\theta(g\mid z)$ and the role-transition $T$ are
identified from the triple laws \eqref{eq:triple} up to a permutation of roles; the permutation is fixed by
semantic \emph{anchors} (top-cosine passages, almost surely $z{=}\textsf{direct}$). This replaces the earlier
neighbour-intuition with tensor-rank machinery and makes ``structure substitutes for replication'' a precise
statement. The Fisher-information analysis of when $\alpha_2-\alpha_1$ is estimable, and the finite-sample rate,
are the accompanying results.}

\subsection{From chains to general corpus graphs}
The chain is the tractable core; the operational setting is a corpus graph. We carry the identification over
under:
\begin{assumption}[Anchors]\label{as:anchor}
A subset of passages have $\Pr(z_i{=}\textsf{direct})\to1$ under the semantic prior (top-cosine passages),
providing partially-known roles that fix the permutation.
\end{assumption}
\begin{assumption}[Graph-smooth truth, judge-blind instrument]\label{as:smooth}
Roles $z$ follow the graph prior \eqref{eq:prior}; the emission law $\Pi$ is conditionally independent of the
graph given $z$ (the judge does not ``know'' the structure).
\end{assumption}
On a general graph, \eqref{eq:triple} generalizes to the joint law over an anchor and its connected
neighbourhood (a local tree/star), giving the multiple views; a connected component with an anchor plays the
role of one long chain. Optional rater/prompt diversity ($R\ge2$ conditionally-independent judges) is classical
replication, which we can add but do not require.
\PROP{Extending the chain theorem to hidden Markov \emph{random fields} on the corpus graph---and identifying
the operating conditions (component connectivity, anchor density, emission separation) under which $\Pi$ and the
role field are identified---is the general-case theoretical program.}

\begin{intuition}
The chain \emph{is} the replication. A lone judge, run once per passage, still leaves fingerprints of its bias
in the \emph{joint} distribution of grades along a reasoning chain, because the chain's latent roles are
correlated in a known way; three consecutive grades are three looks at the middle role. Anchors then say which
role is which. This is exactly dual to the companion paper: there the graph propagated \emph{belief}; here it
manufactures the repeated views that \emph{identify the instrument}.
\end{intuition}

\paragraph{An honest clarification about ``no gold.''} We use no human relevance labels and no repeated
judgments---but we do use \emph{weak semantic anchors} (top-cosine passages, almost surely direct evidence) as
soft pseudo-labels. Their role is only to \emph{name} the latent states \emph{after} identification, not to
create it: without any anchors the model is identified up to a permutation of roles, and the semantics merely
resolve the permutation. We are careful not to imply the method uses no identifying side information whatsoever.

\section{Inference and Calibration}\label{sec:infer}
\subsection{Posterior}
The target is $\Pr(z,\theta,\alpha,\kappa\mid\{g^{(\rho)}\},A)$ under \eqref{eq:obs}--\eqref{eq:prior}. Because
we adopt the ordered-\emph{probit} emission \eqref{eq:probit}, inference uses \emph{Albert--Chib} latent-normal
data augmentation: introduce the latent propensities $\eta_i$, Gaussian given the role $z_i$ and the role
effects $\alpha$, which renders the $\alpha,\kappa$ updates conjugate; the discrete role field $z$ is updated by
message passing. (Had we chosen an ordered-\emph{logit} emission, Pólya--Gamma augmentation would be the
analogue; we commit to probit for its clean latent-normal structure, and we keep the two consistent.)
\PARTIAL{One caveat we will not paper over: the Potts coupling $\theta$ carries a normalizing constant
$Z(\theta)$ that depends on $\theta$, so $\theta$ is \emph{not} a simple conjugate Gibbs update. For the first
version we therefore work on a \emph{chain or tree}, where exact forward--backward message passing is available
and $Z(\theta)$ is tractable---which also matches the identification theorem of \S\ref{sec:ident}; on general
graphs we fix/tune $\theta$ or use pseudolikelihood / exchange (auxiliary-variable) algorithms. Inference is per
query ($N{=}100$) and embarrassingly parallel across queries.}

\subsection{The deliverable: a calibrated relevance posterior}
The output we care about is the bias-corrected posterior $\Pr(r_i{=}1\mid g,A)$---what we should believe about
true relevance after accounting for the judge's structured bias.
\PROP{\textbf{Calibration target.} This posterior should be \emph{calibrated} (its credible sets attain nominal
coverage of the true relevant set), whereas the raw judge is not; standard reliability diagrams / coverage
checks against held-out gold quantify the improvement.}

\subsection{What makes Paper B stand on its own}
A measurement model earns its place on \emph{intrinsic} statistical grounds, before any downstream use. The
primary deliverables are: (i) \emph{identifiable estimation} of the judge's bias---the emission law, and in
particular the bridge-blindness gap $\alpha_2-\alpha_1$; (ii) recovery of latent \emph{role prevalence}; (iii) a
\emph{calibrated} posterior relevance $\Pr(r_i{=}1\mid g,A)$ with correct coverage, and specifically correct
uncertainty on \emph{bridge} items, where the raw judge is confidently wrong; and (iv) improved prediction of
held-out gold relevance versus the raw judge. Demonstrating these makes the paper a contribution to measurement
under biased oracles regardless of the application.
\paragraph{The consequential demonstration.} \emph{Then}, as the payoff, plugging the corrected posterior into
the companion paper's set-completion objective (replacing the raw label in the GP likelihood) should reduce
expected retrieval loss.
\PROP{We will run the retrieval comparison, but not as the \emph{sole} criterion---otherwise a reviewer could
fairly ask why the measurement model is not merely a section of the retrieval paper. The go/no-go is the
\emph{ensemble} (i)--(iv) together with the downstream gain, all of it cheap, since we already hold the judge
grades, gold labels, and graphs.}

\section{Empirical Grounding: What We Already See}\label{sec:evid}
\TESTED{The bias exists and is structured (repeated from \S\ref{sec:model}): naive judge gold-recall $0.35$;
hop-aware prompt $0.82$ (an $\alpha_1$ shift); precision $0.17$--$0.36$; degradation on longer chains and
adversarial corpora. Any adequate model must reproduce these.}
\PARTIAL{Per corpus we hold the joint distribution of judge grades, gold labels, and graph structure over tens
of thousands of (query, passage) pairs. This is enough to \emph{fit} $\Pi$ against a gold-derived role proxy
(bridge $\approx$ semantically buried but graph-connected gold; direct $\approx$ high-cosine gold) and test
role-dependence directly---the first, cheap experiment (``Dig~1''), which grounds the observation law
\eqref{eq:obs}--\eqref{eq:probit} before any of the harder machinery.}

\section{Related Work}
Our proof machinery draws on \emph{measurement-error and misclassification} models (Carroll et al.),
\emph{latent-class rater} models and Bayesian crowdsourcing (Dawid--Skene and successors), \emph{Potts/GMRF}
priors, and \emph{tensor / three-view identifiability} of latent-variable models (Kruskal; Allman, Matias \&
Rhodes) with classical \emph{HMM identifiability} (Gassiat et al.). Two nearby lines demand explicit care.
\emph{First}, a fast-growing 2026 literature already frames the \emph{LLM judge as a measurement instrument}
(noisy-surrogate estimation with a gold calibration sample; IRT-based judge diagnostics; total-evaluation-error
decompositions)---so we claim \emph{no} novelty for that framing; our contribution begins after it.
\emph{Second}, and most pointedly, the \emph{clusterability} line (Zhu et al., 2021; Liu et al., 2023) already
shows that dependence can substitute for repeated labels in identifying a noise-transition matrix---but through
\emph{homophilic}, same-class local structure. Our defensible novelty is the complement: (i) misclassification
that is \emph{differential} in a latent role refining the target; (ii) \emph{heterophilic}, role-\emph{compatible}
graph structure (bridge $\leftrightarrow$ direct), where neighbours \emph{differ}; (iii) single-judge
\emph{structural identification} through that heterophilic dependence, with semantics only naming states; and
(iv) evaluation on intrinsic calibration and a downstream decision rather than agreement rates.

\section{Discussion}
The theme---\emph{rigorous Bayesian treatment of imperfect LLM oracles}---extends well beyond retrieval. The
same model shape covers LLM correctness verifiers (a companion line on text-to-SQL, where a cheap verifier
grades program correctness), automated evaluation of generated text, and weak supervision, wherever a cheap
biased oracle scores structured items. The transportable idea is \emph{structure-as-identification}: when you
can neither replicate the oracle nor collect gold, the dependence structure among the items can de-bias the
instrument that scores them. That is a statistician's contribution to a problem the machine-learning community
currently handles by hand.

\end{document}
